\documentclass[letterpaper,11pt]{article}
\usepackage[top=0.8in, bottom=0.8in, left=0.9in, right=0.9in]{geometry}
\usepackage[hidelinks]{hyperref}
\usepackage{lmodern}
\usepackage{parskip}
\usepackage{microtype}
\setlength{\parskip}{6pt}
\setlength{\parindent}{0pt}
\begin{document}
\begin{flushleft}
\textbf{\large Ajay Venkatesh} \\
Brooklyn, NY \\
+1 516 585 9013 \\
\href{mailto:av3855@nyu.edu}{av3855@nyu.edu}
\end{flushleft}
\vspace{10pt}
May 21, 2026
\vspace{6pt}

Hiring Manager \\
Leadership Circle

\vspace{10pt}

Dear Hiring Manager,

My name is Ajay Venkatesh. I'm finishing my M.S. in Computer Engineering at NYU, with production software experience at PwC in Bengaluru, a summer SDE internship at Amazon, and ongoing on-campus work at NYU's Novel AI Technologies. I'm writing about the Full Stack Software Engineer role supporting Project Center and LeadHub.

The strongest match is the platform work I've owned at NYU. I shipped a customer-facing insurance analytics platform on \textbf{React}, \textbf{Node.js}, and \textbf{REST APIs} backed by \textbf{PostgreSQL}, deployed via \textbf{Docker} on \textbf{EC2}, with \textbf{Auth0-style authentication}, \textbf{RBAC}, row-level security, and Stripe billing for paying customers. I designed the data model end-to-end, wrote and tuned the schema as the product evolved, and supported production users directly. That blend of backend services, frontend experience, and live customer support is exactly what Project Center and LeadHub sound like.

In parallel I shipped \textbf{Node.js microservices} at Campus Mesh powering real-time academic workflows, with a recommendation engine and \textbf{LLM-based retrieval pipelines (RAG)} for an in-app AI assistant, partitioned data modeling, and idempotent async processing. At Amazon I built a \textbf{React 18} frontend with optimized API integration, advanced filtering, and client-side caching, alongside backend services in Java with low-level design patterns and unit / integration testing.

On overlap with your stack: I'm fluent in \textbf{JavaScript / TypeScript}, \textbf{React}, \textbf{Next.js}, \textbf{Node.js}, \textbf{REST API design}, \textbf{PostgreSQL} (query writing, performance tuning, schema evolution), \textbf{Auth0} and \textbf{OAuth} flows, and \textbf{AWS} (\textbf{S3}, \textbf{Lambda}, \textbf{SQS}). I write tests with \textbf{Jest} on the frontend side and treat code review as part of daily work. \textbf{Ruby on Rails} is the one gap, and I'd plan to come up the curve on it quickly inside the existing Project Center codebase while contributing in earnest on the Next.js, NestJS, and Aurora PostgreSQL side from day one.

AI-assisted development is genuinely how I work day to day. I use \textbf{Claude Code} and \textbf{Cursor} for refactors, debugging, code generation, and reasoning across unfamiliar parts of a codebase, and I've thought a lot about where these tools earn their keep versus where they slow you down. Contributing to a team that's actively shaping how AI tooling is woven into delivery is something I'd be excited to do.

What pulls me toward Leadership Circle is the mission. Coaching and leadership development is a domain where the software has to be patient with people, careful with sensitive information, and continuous over months and years rather than minutes. That's a different kind of design constraint than ad tech or e-commerce, and it's the kind of work that compounds into real change in how people show up at work.

I'd welcome the chance to discuss further. Thanks for considering my application.

\vspace{12pt}
Sincerely, \\
Ajay Venkatesh

\end{document}
